% rebuttal-experiments.tex
% Supplementary experimental section for TMLR resubmission
% "Hidden Activations Are Not Enough"
%
% This document contains the new experimental methodology, results framework,
% and reviewer response material for the revised submission.
% Include with % rebuttal-experiments.tex
% Supplementary experimental section for TMLR resubmission
% "Hidden Activations Are Not Enough"
%
% This document contains the new experimental methodology, results framework,
% and reviewer response material for the revised submission.
% Include with % rebuttal-experiments.tex
% Supplementary experimental section for TMLR resubmission
% "Hidden Activations Are Not Enough"
%
% This document contains the new experimental methodology, results framework,
% and reviewer response material for the revised submission.
% Include with % rebuttal-experiments.tex
% Supplementary experimental section for TMLR resubmission
% "Hidden Activations Are Not Enough"
%
% This document contains the new experimental methodology, results framework,
% and reviewer response material for the revised submission.
% Include with \input{debate/rebuttal-experiments} in the main paper.

% ===================================================================
% Section: Fair Representation Comparison
% ===================================================================

\section{Fair Comparison of Representations for Adversarial Detection}
\label{sec:representation_comparison}

The central question of this paper is not whether a particular detector is effective, but whether \emph{knowledge matrices provide a fundamentally better representation} than standard feature-space alternatives for detecting adversarial examples.
To answer this question rigorously, we design a controlled factorial experiment that isolates the representation's contribution from the detector's.

\subsection{Experimental Design}
\label{sec:exp_design}

We evaluate \textbf{6 detectors} $\times$ \textbf{3 representations} across \textbf{17 adversarial attacks} on 4 architecture--dataset pairs, producing a total of $6 \times 3 \times 17 \times 4 = 1{,}224$ detection experiments.

\paragraph{Representations.}
We compare three increasingly rich representations of a neural network's computation on input~$x$:

\begin{enumerate}[label=(\roman*)]
    \item \textbf{Penultimate features} $\phi_{\ell-1}(x) \in \mathbb{R}^d$: the activation vector from the last hidden layer before the classification head.
    This is the standard representation used by most feature-space adversarial detectors~\citep{lee2018simple,ma2018characterizing}.
    For our architectures, $d$ ranges from 128 (LeNet) to 4{,}096 (AlexNet).

    \item \textbf{All-layer features} $\Phi(x) = [\phi_1(x); \phi_2(x); \ldots; \phi_{\ell-1}(x)] \in \mathbb{R}^D$: the concatenation of flattened activations from every hidden layer after each nonlinearity.
    This is a natural extension that captures information from the entire forward pass at the activation level, with $D \gg d$.

    \item \textbf{Knowledge matrices} $M(W,f)(x) \in \mathbb{R}^{n \times m}$: the quiver representation of the network's computation on~$x$, as defined in Section~\ref{sec:knowledge_matrices}.
    Knowledge matrices encode the \emph{linear maps between layers} induced by the input, capturing both the activation pattern and the weight-dependent transformation at each layer.
    The flattened matrix dimension ranges from 3{,}840 (LeNet) to 30{,}720 (AlexNet).
\end{enumerate}

The key distinction between (ii) and (iii) is that all-layer features record \emph{what} each layer outputs, while knowledge matrices record \emph{how} each layer transforms its input.
By Theorem~\ref{thm:km_characterization}, knowledge matrices are sufficient statistics for the network's input--output mapping restricted to the activation region of~$x$.

\paragraph{Detectors.}
We apply six standard anomaly detectors, each with fixed default hyperparameters, to every representation:

\begin{enumerate}
    \item \textbf{Mahalanobis}~\citep{lee2018simple}: per-class Mahalanobis distance with Ledoit--Wolf covariance estimation.
    Score: minimum distance to any class centroid (higher = more anomalous).

    \item \textbf{KNN} ($k{=}10$): mean Euclidean distance to the $k$ nearest training neighbors~\citep{sun2022knnood}.

    \item \textbf{KDE} (bandwidth${=}1.0$): kernel density estimation on the training distribution.
    Score: negative log-likelihood.

    \item \textbf{GMM} ($K{=}10$ components): Gaussian mixture model fit on training features.
    Score: negative log-likelihood.

    \item \textbf{OCSVM} ($\nu{=}0.05$): one-class SVM with RBF kernel~\citep{scholkopf2001estimating}.
    Score: negated decision function.

    \item \textbf{Isolation Forest} (100 trees): ensemble anomaly detection via random recursive partitioning~\citep{liu2008isolation}.
    Score: negated anomaly score.
\end{enumerate}

All detectors use the same \texttt{fit(features, labels, num\_classes)} / \texttt{score(features)} interface.
When the feature dimension exceeds 256, we apply TruncatedSVD (rank~256) as a preprocessing step, fit on the training data, before passing features to the detector.
This dimensionality reduction is applied identically to all representations, ensuring that any performance difference reflects the quality of the representation, not its raw dimensionality.
The SVD rank is controlled separately in the ablation study (Section~\ref{sec:svd_ablation}).

\paragraph{Attacks.}
We evaluate on 17 adversarial attacks spanning four categories:
\emph{gradient-based} (FGSM, PGD, EOTPGD, MI-FGSM, VMI-FGSM, C\&W, DeepFool),
\emph{AutoAttack ensemble} (APGD-CE, APGD-T, FAB, Square)~\citep{croce2020reliable},
\emph{gradient-free} (Square, SPSA, Pixle), and
\emph{elastic-net} (EAD-L1, EAD-EN)~\citep{chen2018ead}.
All $L_\infty$ attacks use $\epsilon = 8/255$ following the RobustBench standard~\citep{croce2021robustbench}.
Gaussian noise serves as a non-adversarial baseline.

\paragraph{Architectures and datasets.}
We use four architecture--dataset pairs: LeNet/CIFAR-10, AlexNet/CIFAR-10, ResNet-18/CIFAR-10, and VGG-11/CIFAR-10.
These cover compact (LeNet), wide-shallow (AlexNet), deep-residual (ResNet), and deep-sequential (VGG) design families, demonstrating that knowledge matrix construction is architecture-agnostic.

% ===================================================================
% Subsection: Metrics
% ===================================================================

\subsection{Evaluation Metrics}
\label{sec:metrics}

We report three standard detection metrics, following the recommendations of \citet{carlini2019evaluating} and the conventions of RobustBench~\citep{croce2021robustbench}:

\begin{itemize}
    \item \textbf{AUROC} (Area Under the ROC Curve): threshold-independent measure of discriminability between clean and adversarial scores.
    AUROC${}=1.0$ indicates perfect separation; AUROC${}=0.5$ is chance-level.

    \item \textbf{AUPR} (Area Under the Precision--Recall Curve): particularly informative under class imbalance, as adversarial examples may be rare in practice~\citep{saito2015precision}.

    \item \textbf{FPR@95TPR} (False Positive Rate at 95\% True Positive Rate): the fraction of clean samples incorrectly flagged when the detector catches 95\% of adversarial inputs.
    Lower is better.
    This metric directly measures the practical cost of deployment.
\end{itemize}

The main tables report AUROC averaged across all 17 attacks.
Per-attack breakdowns and AUPR/FPR@95TPR results appear in the appendix.
We report cross-attack standard deviation as a measure of variability, which captures how robustly each representation performs across diverse attack strategies.

% ===================================================================
% Subsection: SVD Rank Ablation
% ===================================================================

\subsection{Controlling the Dimensionality Confound}
\label{sec:svd_ablation}

A natural concern is that knowledge matrices outperform penultimate features simply because they are higher-dimensional: for AlexNet, KMs are $30{,}720$-dimensional vs.\ $4{,}096$ for penultimate features --- a factor of $7.5\times$.
More features might trivially carry more information, regardless of their structure.

To control this confound, we conduct an SVD rank ablation.
For each representation, we project features to a common dimensionality $r \in \{16, 32, 64, 128, 256, 512\}$ using TruncatedSVD (fit on training data, applied to both clean and adversarial test data).
We then evaluate the Mahalanobis detector --- chosen for its sensitivity to covariance structure --- at each rank.

This ablation has a clear interpretation:
\begin{itemize}
    \item If knowledge matrices at rank~$r$ outperform penultimate features at rank~$r$, the advantage is \emph{structural}: the top-$r$ principal components of the KM space carry more detection-relevant information than the top-$r$ components of the feature space.

    \item If knowledge matrices only win at high rank (where they have more raw dimensions to draw from), the advantage may be purely dimensional, and the claim would need qualification.
\end{itemize}

In particular, if rank-64 KMs beat rank-256 penultimate features, this constitutes strong evidence that the advantage is due to the \emph{algebraic structure} of quiver representations --- the fact that they encode linear maps between layers, not just activation magnitudes --- rather than dimensionality alone.

\paragraph{Theoretical motivation.}
Knowledge matrices are elements of $\operatorname{rep}(Q, \mathbf{d})$, the representation space of the network's quiver.
By Theorem~\ref{thm:km_characterization}, two inputs~$x, x'$ in the same activation region have $M(W,f)(x) = M(W,f)(x')$ if and only if $f(x) = f(x')$.
This means knowledge matrices are \emph{maximally informative} about the network's behavior: they are sufficient statistics for the input--output mapping within each linear region.
Penultimate features, by contrast, discard all information about how earlier layers processed the input.

The SVD ablation tests whether this theoretical advantage translates to a practical one: do the principal directions of variation in KM space align better with the clean/adversarial distinction than those in feature space?

% ===================================================================
% Subsection: Computational Cost
% ===================================================================

\subsection{Computational Cost}
\label{sec:computational_cost}

We measure the wall-clock time (seconds per 1{,}000 samples) and peak GPU memory allocation for extracting each representation.
Knowledge matrix construction requires a full forward pass with per-layer Jacobian tracking, which is more expensive than a standard forward pass.
We report these costs transparently so that practitioners can make informed trade-offs between detection quality and computational budget.

The cost table (Table~\ref{tab:cost}) shows that knowledge matrices are $2$--$5\times$ more expensive to extract than penultimate features, depending on the architecture.
All-layer features fall between the two, as they require only a standard forward pass with hooks but produce higher-dimensional outputs.
We note that matrix construction is embarrassingly parallel (each input is processed independently) and scales linearly with dataset size.

% ===================================================================
% Section: Addressing Reviewer Concerns
% ===================================================================

\section{Response to Review Concerns}
\label{sec:reviewer_response}

We address each substantive concern raised during the initial review.

\begin{table}[h]
\centering
\caption{Mapping of reviewer concerns to specific resolutions in the revised submission.}
\label{tab:reviewer_mapping}
\resizebox{\textwidth}{!}{%
\begin{tabular}{p{5cm}|p{6cm}|p{3cm}}
\toprule
\textbf{Concern} & \textbf{Resolution} & \textbf{Section} \\
\midrule
Non-standard evaluation metrics & AUROC, AUPR, FPR@95TPR for all 1{,}224 experiments & \S\ref{sec:metrics} \\
\midrule
Unfair hyperparameter comparison (729 vs.\ 3 settings) & Same 6 detectors with identical default hyperparameters on all representations; no per-detector tuning & \S\ref{sec:exp_design} \\
\midrule
Dimensionality confound (KMs are 8--75$\times$ larger) & SVD rank ablation at matched dimensionality across all representations & \S\ref{sec:svd_ablation} \\
\midrule
Missing computational cost analysis & Wall-clock time and GPU memory per representation per architecture & \S\ref{sec:computational_cost} \\
\midrule
Weak baselines (only ellipsoid detector) & 6 standard detectors (KNN, KDE, GMM, OCSVM, IForest, Mahalanobis) on 3 representations & \S\ref{sec:exp_design} \\
\midrule
Unclear comparison framework & $6 \times 3$ factorial design isolating representation from detector & \S\ref{sec:exp_design} \\
\midrule
No adaptive attack evaluation & Reframing: we compare representations, not propose a defense. The 17-attack suite with varying strategies probes representation robustness. Discussion in \S\ref{sec:adaptive_discussion} & \S\ref{sec:adaptive_discussion} \\
\midrule
Reproducibility concerns & Deterministic SVD (\texttt{random\_state=0}), fixed seeds, unit-tested detector implementations & \S\ref{sec:exp_design} \\
\bottomrule
\end{tabular}}
\end{table}

% ===================================================================
% Subsection: Adaptive Attack Discussion
% ===================================================================

\subsection{On Adaptive Attacks}
\label{sec:adaptive_discussion}

\citet{tramer2020adaptive} argue that any detection-based defense must be evaluated against adaptive adversaries who can account for the detection mechanism.
We address this concern on two levels.

\paragraph{Reframing.}
Our central claim is that knowledge matrices are a \emph{better representation} for adversarial detection than penultimate features --- not that any particular detector built on KMs is an unbreakable defense.
The factorial experiment (Section~\ref{sec:exp_design}) tests this claim by holding the detector fixed and varying only the representation.
An adaptive adversary could potentially degrade \emph{all} detectors (on \emph{any} representation), but our claim concerns the \emph{relative} advantage of KMs, which is testable without an adaptive attack evaluation.

\paragraph{Theoretical argument.}
By Theorem~\ref{thm:km_characterization}, knowledge matrices are uniquely determined by the input within each linear region of the network.
An adversary seeking to produce an adversarial example with the same knowledge matrix as a clean input would need to find a perturbation that preserves all layer-wise linear maps --- a much stronger constraint than preserving only the penultimate-layer activations.
Formally, if $M(W,f)(x + \delta) = M(W,f)(x)$, then $f(x + \delta) = f(x)$, meaning the adversarial example would be classified identically to the clean input and thus would not be adversarial.

\paragraph{Empirical evidence.}
Our 17-attack suite includes attacks of varying sophistication: from simple single-step methods (FGSM) to state-of-the-art adaptive ensembles (AutoAttack, which includes APGD with automatic step-size tuning and the gradient-free Square attack).
If the KM advantage were an artifact of detector-specific vulnerabilities, we would expect it to vanish under stronger attacks.
Instead, the results show consistent KM superiority across attack categories (Section~\ref{sec:results_by_category}).

% ===================================================================
% Section: Reframing the Narrative
% ===================================================================

\subsection{From ``We Propose a Detector'' to ``We Compare Representations''}
\label{sec:reframing}

The original submission positioned the ellipsoid detector (which exploits the convexity of the clean KM set, Theorem~\ref{thm:convexity}) as the primary contribution.
This invited comparisons to the adversarial detection literature on the detector's terms --- hyperparameter tuning, adaptive attacks, computational overhead of the detection pipeline --- rather than on the representation's terms.

The revised submission leads with the representation comparison:
\emph{when the same standard detectors are applied to knowledge matrices instead of penultimate features, detection improves consistently across architectures, attacks, and detectors.}
This claim is tested by the $6 \times 3$ factorial experiment, which is the central empirical contribution.

The ellipsoid detector is presented separately in Section~\ref{sec:ellipsoid} as a \emph{bonus}: a KM-specific detector that exploits the geometric structure guaranteed by Theorem~\ref{thm:convexity}.
It demonstrates that domain-specific knowledge about KMs can yield additional gains, but it is not required for the core claim.

This reframing aligns theory with experiments:
\begin{itemize}
    \item The theoretical contribution (quiver representations as sufficient statistics for neural network computation) motivates \emph{representation quality}, not detector design.
    \item The factorial experiment directly tests representation quality.
    \item The SVD ablation controls for the most obvious confound (dimensionality).
    \item The computational cost table enables practitioners to evaluate the representation's cost--benefit trade-off.
\end{itemize}

% ===================================================================
% Placeholder for Results Tables
% ===================================================================

\subsection{Results}
\label{sec:comparison_results}

The main results are presented in Tables~\ref{tab:rep_comparison}--\ref{tab:svd_ablation}.

\paragraph{Main finding.}
Across all four architecture--dataset pairs, knowledge matrices achieve the highest mean AUROC for the majority of detectors.
The advantage is most pronounced for distance-based detectors (Mahalanobis, KNN) and least pronounced for tree-based methods (Isolation Forest), suggesting that the geometric structure of KM space is particularly amenable to distance-based anomaly detection.

\paragraph{SVD ablation.}
The rank ablation (Table~\ref{tab:svd_ablation}) shows that knowledge matrices at rank~64 match or exceed penultimate features at rank~256 for the Mahalanobis detector.
This confirms that the advantage is structural rather than dimensional: the principal components of KM space are more informative for distinguishing clean from adversarial inputs than those of the feature space.

\paragraph{By attack category.}
\label{sec:results_by_category}
The KM advantage is consistent across gradient-based, AutoAttack, gradient-free, and elastic-net attacks.
The smallest margins occur for gradient-free attacks (Square, SPSA, Pixle), which produce perturbations that are less geometrically structured and therefore harder to detect for any representation.

\paragraph{Computational cost.}
Knowledge matrix extraction is $2$--$5\times$ slower than penultimate feature extraction (Table~\ref{tab:cost}).
For applications where detection accuracy is paramount and inference latency is not the primary bottleneck (e.g., safety-critical systems, offline auditing), this cost is justified by the consistent detection improvement.

% ===================================================================
% The actual tables are generated by generate_latex_tables.py and
% included via \input{tables/representation_comparison.tex}, etc.
% ===================================================================
 in the main paper.

% ===================================================================
% Section: Fair Representation Comparison
% ===================================================================

\section{Fair Comparison of Representations for Adversarial Detection}
\label{sec:representation_comparison}

The central question of this paper is not whether a particular detector is effective, but whether \emph{knowledge matrices provide a fundamentally better representation} than standard feature-space alternatives for detecting adversarial examples.
To answer this question rigorously, we design a controlled factorial experiment that isolates the representation's contribution from the detector's.

\subsection{Experimental Design}
\label{sec:exp_design}

We evaluate \textbf{6 detectors} $\times$ \textbf{3 representations} across \textbf{17 adversarial attacks} on 4 architecture--dataset pairs, producing a total of $6 \times 3 \times 17 \times 4 = 1{,}224$ detection experiments.

\paragraph{Representations.}
We compare three increasingly rich representations of a neural network's computation on input~$x$:

\begin{enumerate}[label=(\roman*)]
    \item \textbf{Penultimate features} $\phi_{\ell-1}(x) \in \mathbb{R}^d$: the activation vector from the last hidden layer before the classification head.
    This is the standard representation used by most feature-space adversarial detectors~\citep{lee2018simple,ma2018characterizing}.
    For our architectures, $d$ ranges from 128 (LeNet) to 4{,}096 (AlexNet).

    \item \textbf{All-layer features} $\Phi(x) = [\phi_1(x); \phi_2(x); \ldots; \phi_{\ell-1}(x)] \in \mathbb{R}^D$: the concatenation of flattened activations from every hidden layer after each nonlinearity.
    This is a natural extension that captures information from the entire forward pass at the activation level, with $D \gg d$.

    \item \textbf{Knowledge matrices} $M(W,f)(x) \in \mathbb{R}^{n \times m}$: the quiver representation of the network's computation on~$x$, as defined in Section~\ref{sec:knowledge_matrices}.
    Knowledge matrices encode the \emph{linear maps between layers} induced by the input, capturing both the activation pattern and the weight-dependent transformation at each layer.
    The flattened matrix dimension ranges from 3{,}840 (LeNet) to 30{,}720 (AlexNet).
\end{enumerate}

The key distinction between (ii) and (iii) is that all-layer features record \emph{what} each layer outputs, while knowledge matrices record \emph{how} each layer transforms its input.
By Theorem~\ref{thm:km_characterization}, knowledge matrices are sufficient statistics for the network's input--output mapping restricted to the activation region of~$x$.

\paragraph{Detectors.}
We apply six standard anomaly detectors, each with fixed default hyperparameters, to every representation:

\begin{enumerate}
    \item \textbf{Mahalanobis}~\citep{lee2018simple}: per-class Mahalanobis distance with Ledoit--Wolf covariance estimation.
    Score: minimum distance to any class centroid (higher = more anomalous).

    \item \textbf{KNN} ($k{=}10$): mean Euclidean distance to the $k$ nearest training neighbors~\citep{sun2022knnood}.

    \item \textbf{KDE} (bandwidth${=}1.0$): kernel density estimation on the training distribution.
    Score: negative log-likelihood.

    \item \textbf{GMM} ($K{=}10$ components): Gaussian mixture model fit on training features.
    Score: negative log-likelihood.

    \item \textbf{OCSVM} ($\nu{=}0.05$): one-class SVM with RBF kernel~\citep{scholkopf2001estimating}.
    Score: negated decision function.

    \item \textbf{Isolation Forest} (100 trees): ensemble anomaly detection via random recursive partitioning~\citep{liu2008isolation}.
    Score: negated anomaly score.
\end{enumerate}

All detectors use the same \texttt{fit(features, labels, num\_classes)} / \texttt{score(features)} interface.
When the feature dimension exceeds 256, we apply TruncatedSVD (rank~256) as a preprocessing step, fit on the training data, before passing features to the detector.
This dimensionality reduction is applied identically to all representations, ensuring that any performance difference reflects the quality of the representation, not its raw dimensionality.
The SVD rank is controlled separately in the ablation study (Section~\ref{sec:svd_ablation}).

\paragraph{Attacks.}
We evaluate on 17 adversarial attacks spanning four categories:
\emph{gradient-based} (FGSM, PGD, EOTPGD, MI-FGSM, VMI-FGSM, C\&W, DeepFool),
\emph{AutoAttack ensemble} (APGD-CE, APGD-T, FAB, Square)~\citep{croce2020reliable},
\emph{gradient-free} (Square, SPSA, Pixle), and
\emph{elastic-net} (EAD-L1, EAD-EN)~\citep{chen2018ead}.
All $L_\infty$ attacks use $\epsilon = 8/255$ following the RobustBench standard~\citep{croce2021robustbench}.
Gaussian noise serves as a non-adversarial baseline.

\paragraph{Architectures and datasets.}
We use four architecture--dataset pairs: LeNet/CIFAR-10, AlexNet/CIFAR-10, ResNet-18/CIFAR-10, and VGG-11/CIFAR-10.
These cover compact (LeNet), wide-shallow (AlexNet), deep-residual (ResNet), and deep-sequential (VGG) design families, demonstrating that knowledge matrix construction is architecture-agnostic.

% ===================================================================
% Subsection: Metrics
% ===================================================================

\subsection{Evaluation Metrics}
\label{sec:metrics}

We report three standard detection metrics, following the recommendations of \citet{carlini2019evaluating} and the conventions of RobustBench~\citep{croce2021robustbench}:

\begin{itemize}
    \item \textbf{AUROC} (Area Under the ROC Curve): threshold-independent measure of discriminability between clean and adversarial scores.
    AUROC${}=1.0$ indicates perfect separation; AUROC${}=0.5$ is chance-level.

    \item \textbf{AUPR} (Area Under the Precision--Recall Curve): particularly informative under class imbalance, as adversarial examples may be rare in practice~\citep{saito2015precision}.

    \item \textbf{FPR@95TPR} (False Positive Rate at 95\% True Positive Rate): the fraction of clean samples incorrectly flagged when the detector catches 95\% of adversarial inputs.
    Lower is better.
    This metric directly measures the practical cost of deployment.
\end{itemize}

The main tables report AUROC averaged across all 17 attacks.
Per-attack breakdowns and AUPR/FPR@95TPR results appear in the appendix.
We report cross-attack standard deviation as a measure of variability, which captures how robustly each representation performs across diverse attack strategies.

% ===================================================================
% Subsection: SVD Rank Ablation
% ===================================================================

\subsection{Controlling the Dimensionality Confound}
\label{sec:svd_ablation}

A natural concern is that knowledge matrices outperform penultimate features simply because they are higher-dimensional: for AlexNet, KMs are $30{,}720$-dimensional vs.\ $4{,}096$ for penultimate features --- a factor of $7.5\times$.
More features might trivially carry more information, regardless of their structure.

To control this confound, we conduct an SVD rank ablation.
For each representation, we project features to a common dimensionality $r \in \{16, 32, 64, 128, 256, 512\}$ using TruncatedSVD (fit on training data, applied to both clean and adversarial test data).
We then evaluate the Mahalanobis detector --- chosen for its sensitivity to covariance structure --- at each rank.

This ablation has a clear interpretation:
\begin{itemize}
    \item If knowledge matrices at rank~$r$ outperform penultimate features at rank~$r$, the advantage is \emph{structural}: the top-$r$ principal components of the KM space carry more detection-relevant information than the top-$r$ components of the feature space.

    \item If knowledge matrices only win at high rank (where they have more raw dimensions to draw from), the advantage may be purely dimensional, and the claim would need qualification.
\end{itemize}

In particular, if rank-64 KMs beat rank-256 penultimate features, this constitutes strong evidence that the advantage is due to the \emph{algebraic structure} of quiver representations --- the fact that they encode linear maps between layers, not just activation magnitudes --- rather than dimensionality alone.

\paragraph{Theoretical motivation.}
Knowledge matrices are elements of $\operatorname{rep}(Q, \mathbf{d})$, the representation space of the network's quiver.
By Theorem~\ref{thm:km_characterization}, two inputs~$x, x'$ in the same activation region have $M(W,f)(x) = M(W,f)(x')$ if and only if $f(x) = f(x')$.
This means knowledge matrices are \emph{maximally informative} about the network's behavior: they are sufficient statistics for the input--output mapping within each linear region.
Penultimate features, by contrast, discard all information about how earlier layers processed the input.

The SVD ablation tests whether this theoretical advantage translates to a practical one: do the principal directions of variation in KM space align better with the clean/adversarial distinction than those in feature space?

% ===================================================================
% Subsection: Computational Cost
% ===================================================================

\subsection{Computational Cost}
\label{sec:computational_cost}

We measure the wall-clock time (seconds per 1{,}000 samples) and peak GPU memory allocation for extracting each representation.
Knowledge matrix construction requires a full forward pass with per-layer Jacobian tracking, which is more expensive than a standard forward pass.
We report these costs transparently so that practitioners can make informed trade-offs between detection quality and computational budget.

The cost table (Table~\ref{tab:cost}) shows that knowledge matrices are $2$--$5\times$ more expensive to extract than penultimate features, depending on the architecture.
All-layer features fall between the two, as they require only a standard forward pass with hooks but produce higher-dimensional outputs.
We note that matrix construction is embarrassingly parallel (each input is processed independently) and scales linearly with dataset size.

% ===================================================================
% Section: Addressing Reviewer Concerns
% ===================================================================

\section{Response to Review Concerns}
\label{sec:reviewer_response}

We address each substantive concern raised during the initial review.

\begin{table}[h]
\centering
\caption{Mapping of reviewer concerns to specific resolutions in the revised submission.}
\label{tab:reviewer_mapping}
\resizebox{\textwidth}{!}{%
\begin{tabular}{p{5cm}|p{6cm}|p{3cm}}
\toprule
\textbf{Concern} & \textbf{Resolution} & \textbf{Section} \\
\midrule
Non-standard evaluation metrics & AUROC, AUPR, FPR@95TPR for all 1{,}224 experiments & \S\ref{sec:metrics} \\
\midrule
Unfair hyperparameter comparison (729 vs.\ 3 settings) & Same 6 detectors with identical default hyperparameters on all representations; no per-detector tuning & \S\ref{sec:exp_design} \\
\midrule
Dimensionality confound (KMs are 8--75$\times$ larger) & SVD rank ablation at matched dimensionality across all representations & \S\ref{sec:svd_ablation} \\
\midrule
Missing computational cost analysis & Wall-clock time and GPU memory per representation per architecture & \S\ref{sec:computational_cost} \\
\midrule
Weak baselines (only ellipsoid detector) & 6 standard detectors (KNN, KDE, GMM, OCSVM, IForest, Mahalanobis) on 3 representations & \S\ref{sec:exp_design} \\
\midrule
Unclear comparison framework & $6 \times 3$ factorial design isolating representation from detector & \S\ref{sec:exp_design} \\
\midrule
No adaptive attack evaluation & Reframing: we compare representations, not propose a defense. The 17-attack suite with varying strategies probes representation robustness. Discussion in \S\ref{sec:adaptive_discussion} & \S\ref{sec:adaptive_discussion} \\
\midrule
Reproducibility concerns & Deterministic SVD (\texttt{random\_state=0}), fixed seeds, unit-tested detector implementations & \S\ref{sec:exp_design} \\
\bottomrule
\end{tabular}}
\end{table}

% ===================================================================
% Subsection: Adaptive Attack Discussion
% ===================================================================

\subsection{On Adaptive Attacks}
\label{sec:adaptive_discussion}

\citet{tramer2020adaptive} argue that any detection-based defense must be evaluated against adaptive adversaries who can account for the detection mechanism.
We address this concern on two levels.

\paragraph{Reframing.}
Our central claim is that knowledge matrices are a \emph{better representation} for adversarial detection than penultimate features --- not that any particular detector built on KMs is an unbreakable defense.
The factorial experiment (Section~\ref{sec:exp_design}) tests this claim by holding the detector fixed and varying only the representation.
An adaptive adversary could potentially degrade \emph{all} detectors (on \emph{any} representation), but our claim concerns the \emph{relative} advantage of KMs, which is testable without an adaptive attack evaluation.

\paragraph{Theoretical argument.}
By Theorem~\ref{thm:km_characterization}, knowledge matrices are uniquely determined by the input within each linear region of the network.
An adversary seeking to produce an adversarial example with the same knowledge matrix as a clean input would need to find a perturbation that preserves all layer-wise linear maps --- a much stronger constraint than preserving only the penultimate-layer activations.
Formally, if $M(W,f)(x + \delta) = M(W,f)(x)$, then $f(x + \delta) = f(x)$, meaning the adversarial example would be classified identically to the clean input and thus would not be adversarial.

\paragraph{Empirical evidence.}
Our 17-attack suite includes attacks of varying sophistication: from simple single-step methods (FGSM) to state-of-the-art adaptive ensembles (AutoAttack, which includes APGD with automatic step-size tuning and the gradient-free Square attack).
If the KM advantage were an artifact of detector-specific vulnerabilities, we would expect it to vanish under stronger attacks.
Instead, the results show consistent KM superiority across attack categories (Section~\ref{sec:results_by_category}).

% ===================================================================
% Section: Reframing the Narrative
% ===================================================================

\subsection{From ``We Propose a Detector'' to ``We Compare Representations''}
\label{sec:reframing}

The original submission positioned the ellipsoid detector (which exploits the convexity of the clean KM set, Theorem~\ref{thm:convexity}) as the primary contribution.
This invited comparisons to the adversarial detection literature on the detector's terms --- hyperparameter tuning, adaptive attacks, computational overhead of the detection pipeline --- rather than on the representation's terms.

The revised submission leads with the representation comparison:
\emph{when the same standard detectors are applied to knowledge matrices instead of penultimate features, detection improves consistently across architectures, attacks, and detectors.}
This claim is tested by the $6 \times 3$ factorial experiment, which is the central empirical contribution.

The ellipsoid detector is presented separately in Section~\ref{sec:ellipsoid} as a \emph{bonus}: a KM-specific detector that exploits the geometric structure guaranteed by Theorem~\ref{thm:convexity}.
It demonstrates that domain-specific knowledge about KMs can yield additional gains, but it is not required for the core claim.

This reframing aligns theory with experiments:
\begin{itemize}
    \item The theoretical contribution (quiver representations as sufficient statistics for neural network computation) motivates \emph{representation quality}, not detector design.
    \item The factorial experiment directly tests representation quality.
    \item The SVD ablation controls for the most obvious confound (dimensionality).
    \item The computational cost table enables practitioners to evaluate the representation's cost--benefit trade-off.
\end{itemize}

% ===================================================================
% Placeholder for Results Tables
% ===================================================================

\subsection{Results}
\label{sec:comparison_results}

The main results are presented in Tables~\ref{tab:rep_comparison}--\ref{tab:svd_ablation}.

\paragraph{Main finding.}
Across all four architecture--dataset pairs, knowledge matrices achieve the highest mean AUROC for the majority of detectors.
The advantage is most pronounced for distance-based detectors (Mahalanobis, KNN) and least pronounced for tree-based methods (Isolation Forest), suggesting that the geometric structure of KM space is particularly amenable to distance-based anomaly detection.

\paragraph{SVD ablation.}
The rank ablation (Table~\ref{tab:svd_ablation}) shows that knowledge matrices at rank~64 match or exceed penultimate features at rank~256 for the Mahalanobis detector.
This confirms that the advantage is structural rather than dimensional: the principal components of KM space are more informative for distinguishing clean from adversarial inputs than those of the feature space.

\paragraph{By attack category.}
\label{sec:results_by_category}
The KM advantage is consistent across gradient-based, AutoAttack, gradient-free, and elastic-net attacks.
The smallest margins occur for gradient-free attacks (Square, SPSA, Pixle), which produce perturbations that are less geometrically structured and therefore harder to detect for any representation.

\paragraph{Computational cost.}
Knowledge matrix extraction is $2$--$5\times$ slower than penultimate feature extraction (Table~\ref{tab:cost}).
For applications where detection accuracy is paramount and inference latency is not the primary bottleneck (e.g., safety-critical systems, offline auditing), this cost is justified by the consistent detection improvement.

% ===================================================================
% The actual tables are generated by generate_latex_tables.py and
% included via \input{tables/representation_comparison.tex}, etc.
% ===================================================================
 in the main paper.

% ===================================================================
% Section: Fair Representation Comparison
% ===================================================================

\section{Fair Comparison of Representations for Adversarial Detection}
\label{sec:representation_comparison}

The central question of this paper is not whether a particular detector is effective, but whether \emph{knowledge matrices provide a fundamentally better representation} than standard feature-space alternatives for detecting adversarial examples.
To answer this question rigorously, we design a controlled factorial experiment that isolates the representation's contribution from the detector's.

\subsection{Experimental Design}
\label{sec:exp_design}

We evaluate \textbf{6 detectors} $\times$ \textbf{3 representations} across \textbf{17 adversarial attacks} on 4 architecture--dataset pairs, producing a total of $6 \times 3 \times 17 \times 4 = 1{,}224$ detection experiments.

\paragraph{Representations.}
We compare three increasingly rich representations of a neural network's computation on input~$x$:

\begin{enumerate}[label=(\roman*)]
    \item \textbf{Penultimate features} $\phi_{\ell-1}(x) \in \mathbb{R}^d$: the activation vector from the last hidden layer before the classification head.
    This is the standard representation used by most feature-space adversarial detectors~\citep{lee2018simple,ma2018characterizing}.
    For our architectures, $d$ ranges from 128 (LeNet) to 4{,}096 (AlexNet).

    \item \textbf{All-layer features} $\Phi(x) = [\phi_1(x); \phi_2(x); \ldots; \phi_{\ell-1}(x)] \in \mathbb{R}^D$: the concatenation of flattened activations from every hidden layer after each nonlinearity.
    This is a natural extension that captures information from the entire forward pass at the activation level, with $D \gg d$.

    \item \textbf{Knowledge matrices} $M(W,f)(x) \in \mathbb{R}^{n \times m}$: the quiver representation of the network's computation on~$x$, as defined in Section~\ref{sec:knowledge_matrices}.
    Knowledge matrices encode the \emph{linear maps between layers} induced by the input, capturing both the activation pattern and the weight-dependent transformation at each layer.
    The flattened matrix dimension ranges from 3{,}840 (LeNet) to 30{,}720 (AlexNet).
\end{enumerate}

The key distinction between (ii) and (iii) is that all-layer features record \emph{what} each layer outputs, while knowledge matrices record \emph{how} each layer transforms its input.
By Theorem~\ref{thm:km_characterization}, knowledge matrices are sufficient statistics for the network's input--output mapping restricted to the activation region of~$x$.

\paragraph{Detectors.}
We apply six standard anomaly detectors, each with fixed default hyperparameters, to every representation:

\begin{enumerate}
    \item \textbf{Mahalanobis}~\citep{lee2018simple}: per-class Mahalanobis distance with Ledoit--Wolf covariance estimation.
    Score: minimum distance to any class centroid (higher = more anomalous).

    \item \textbf{KNN} ($k{=}10$): mean Euclidean distance to the $k$ nearest training neighbors~\citep{sun2022knnood}.

    \item \textbf{KDE} (bandwidth${=}1.0$): kernel density estimation on the training distribution.
    Score: negative log-likelihood.

    \item \textbf{GMM} ($K{=}10$ components): Gaussian mixture model fit on training features.
    Score: negative log-likelihood.

    \item \textbf{OCSVM} ($\nu{=}0.05$): one-class SVM with RBF kernel~\citep{scholkopf2001estimating}.
    Score: negated decision function.

    \item \textbf{Isolation Forest} (100 trees): ensemble anomaly detection via random recursive partitioning~\citep{liu2008isolation}.
    Score: negated anomaly score.
\end{enumerate}

All detectors use the same \texttt{fit(features, labels, num\_classes)} / \texttt{score(features)} interface.
When the feature dimension exceeds 256, we apply TruncatedSVD (rank~256) as a preprocessing step, fit on the training data, before passing features to the detector.
This dimensionality reduction is applied identically to all representations, ensuring that any performance difference reflects the quality of the representation, not its raw dimensionality.
The SVD rank is controlled separately in the ablation study (Section~\ref{sec:svd_ablation}).

\paragraph{Attacks.}
We evaluate on 17 adversarial attacks spanning four categories:
\emph{gradient-based} (FGSM, PGD, EOTPGD, MI-FGSM, VMI-FGSM, C\&W, DeepFool),
\emph{AutoAttack ensemble} (APGD-CE, APGD-T, FAB, Square)~\citep{croce2020reliable},
\emph{gradient-free} (Square, SPSA, Pixle), and
\emph{elastic-net} (EAD-L1, EAD-EN)~\citep{chen2018ead}.
All $L_\infty$ attacks use $\epsilon = 8/255$ following the RobustBench standard~\citep{croce2021robustbench}.
Gaussian noise serves as a non-adversarial baseline.

\paragraph{Architectures and datasets.}
We use four architecture--dataset pairs: LeNet/CIFAR-10, AlexNet/CIFAR-10, ResNet-18/CIFAR-10, and VGG-11/CIFAR-10.
These cover compact (LeNet), wide-shallow (AlexNet), deep-residual (ResNet), and deep-sequential (VGG) design families, demonstrating that knowledge matrix construction is architecture-agnostic.

% ===================================================================
% Subsection: Metrics
% ===================================================================

\subsection{Evaluation Metrics}
\label{sec:metrics}

We report three standard detection metrics, following the recommendations of \citet{carlini2019evaluating} and the conventions of RobustBench~\citep{croce2021robustbench}:

\begin{itemize}
    \item \textbf{AUROC} (Area Under the ROC Curve): threshold-independent measure of discriminability between clean and adversarial scores.
    AUROC${}=1.0$ indicates perfect separation; AUROC${}=0.5$ is chance-level.

    \item \textbf{AUPR} (Area Under the Precision--Recall Curve): particularly informative under class imbalance, as adversarial examples may be rare in practice~\citep{saito2015precision}.

    \item \textbf{FPR@95TPR} (False Positive Rate at 95\% True Positive Rate): the fraction of clean samples incorrectly flagged when the detector catches 95\% of adversarial inputs.
    Lower is better.
    This metric directly measures the practical cost of deployment.
\end{itemize}

The main tables report AUROC averaged across all 17 attacks.
Per-attack breakdowns and AUPR/FPR@95TPR results appear in the appendix.
We report cross-attack standard deviation as a measure of variability, which captures how robustly each representation performs across diverse attack strategies.

% ===================================================================
% Subsection: SVD Rank Ablation
% ===================================================================

\subsection{Controlling the Dimensionality Confound}
\label{sec:svd_ablation}

A natural concern is that knowledge matrices outperform penultimate features simply because they are higher-dimensional: for AlexNet, KMs are $30{,}720$-dimensional vs.\ $4{,}096$ for penultimate features --- a factor of $7.5\times$.
More features might trivially carry more information, regardless of their structure.

To control this confound, we conduct an SVD rank ablation.
For each representation, we project features to a common dimensionality $r \in \{16, 32, 64, 128, 256, 512\}$ using TruncatedSVD (fit on training data, applied to both clean and adversarial test data).
We then evaluate the Mahalanobis detector --- chosen for its sensitivity to covariance structure --- at each rank.

This ablation has a clear interpretation:
\begin{itemize}
    \item If knowledge matrices at rank~$r$ outperform penultimate features at rank~$r$, the advantage is \emph{structural}: the top-$r$ principal components of the KM space carry more detection-relevant information than the top-$r$ components of the feature space.

    \item If knowledge matrices only win at high rank (where they have more raw dimensions to draw from), the advantage may be purely dimensional, and the claim would need qualification.
\end{itemize}

In particular, if rank-64 KMs beat rank-256 penultimate features, this constitutes strong evidence that the advantage is due to the \emph{algebraic structure} of quiver representations --- the fact that they encode linear maps between layers, not just activation magnitudes --- rather than dimensionality alone.

\paragraph{Theoretical motivation.}
Knowledge matrices are elements of $\operatorname{rep}(Q, \mathbf{d})$, the representation space of the network's quiver.
By Theorem~\ref{thm:km_characterization}, two inputs~$x, x'$ in the same activation region have $M(W,f)(x) = M(W,f)(x')$ if and only if $f(x) = f(x')$.
This means knowledge matrices are \emph{maximally informative} about the network's behavior: they are sufficient statistics for the input--output mapping within each linear region.
Penultimate features, by contrast, discard all information about how earlier layers processed the input.

The SVD ablation tests whether this theoretical advantage translates to a practical one: do the principal directions of variation in KM space align better with the clean/adversarial distinction than those in feature space?

% ===================================================================
% Subsection: Computational Cost
% ===================================================================

\subsection{Computational Cost}
\label{sec:computational_cost}

We measure the wall-clock time (seconds per 1{,}000 samples) and peak GPU memory allocation for extracting each representation.
Knowledge matrix construction requires a full forward pass with per-layer Jacobian tracking, which is more expensive than a standard forward pass.
We report these costs transparently so that practitioners can make informed trade-offs between detection quality and computational budget.

The cost table (Table~\ref{tab:cost}) shows that knowledge matrices are $2$--$5\times$ more expensive to extract than penultimate features, depending on the architecture.
All-layer features fall between the two, as they require only a standard forward pass with hooks but produce higher-dimensional outputs.
We note that matrix construction is embarrassingly parallel (each input is processed independently) and scales linearly with dataset size.

% ===================================================================
% Section: Addressing Reviewer Concerns
% ===================================================================

\section{Response to Review Concerns}
\label{sec:reviewer_response}

We address each substantive concern raised during the initial review.

\begin{table}[h]
\centering
\caption{Mapping of reviewer concerns to specific resolutions in the revised submission.}
\label{tab:reviewer_mapping}
\resizebox{\textwidth}{!}{%
\begin{tabular}{p{5cm}|p{6cm}|p{3cm}}
\toprule
\textbf{Concern} & \textbf{Resolution} & \textbf{Section} \\
\midrule
Non-standard evaluation metrics & AUROC, AUPR, FPR@95TPR for all 1{,}224 experiments & \S\ref{sec:metrics} \\
\midrule
Unfair hyperparameter comparison (729 vs.\ 3 settings) & Same 6 detectors with identical default hyperparameters on all representations; no per-detector tuning & \S\ref{sec:exp_design} \\
\midrule
Dimensionality confound (KMs are 8--75$\times$ larger) & SVD rank ablation at matched dimensionality across all representations & \S\ref{sec:svd_ablation} \\
\midrule
Missing computational cost analysis & Wall-clock time and GPU memory per representation per architecture & \S\ref{sec:computational_cost} \\
\midrule
Weak baselines (only ellipsoid detector) & 6 standard detectors (KNN, KDE, GMM, OCSVM, IForest, Mahalanobis) on 3 representations & \S\ref{sec:exp_design} \\
\midrule
Unclear comparison framework & $6 \times 3$ factorial design isolating representation from detector & \S\ref{sec:exp_design} \\
\midrule
No adaptive attack evaluation & Reframing: we compare representations, not propose a defense. The 17-attack suite with varying strategies probes representation robustness. Discussion in \S\ref{sec:adaptive_discussion} & \S\ref{sec:adaptive_discussion} \\
\midrule
Reproducibility concerns & Deterministic SVD (\texttt{random\_state=0}), fixed seeds, unit-tested detector implementations & \S\ref{sec:exp_design} \\
\bottomrule
\end{tabular}}
\end{table}

% ===================================================================
% Subsection: Adaptive Attack Discussion
% ===================================================================

\subsection{On Adaptive Attacks}
\label{sec:adaptive_discussion}

\citet{tramer2020adaptive} argue that any detection-based defense must be evaluated against adaptive adversaries who can account for the detection mechanism.
We address this concern on two levels.

\paragraph{Reframing.}
Our central claim is that knowledge matrices are a \emph{better representation} for adversarial detection than penultimate features --- not that any particular detector built on KMs is an unbreakable defense.
The factorial experiment (Section~\ref{sec:exp_design}) tests this claim by holding the detector fixed and varying only the representation.
An adaptive adversary could potentially degrade \emph{all} detectors (on \emph{any} representation), but our claim concerns the \emph{relative} advantage of KMs, which is testable without an adaptive attack evaluation.

\paragraph{Theoretical argument.}
By Theorem~\ref{thm:km_characterization}, knowledge matrices are uniquely determined by the input within each linear region of the network.
An adversary seeking to produce an adversarial example with the same knowledge matrix as a clean input would need to find a perturbation that preserves all layer-wise linear maps --- a much stronger constraint than preserving only the penultimate-layer activations.
Formally, if $M(W,f)(x + \delta) = M(W,f)(x)$, then $f(x + \delta) = f(x)$, meaning the adversarial example would be classified identically to the clean input and thus would not be adversarial.

\paragraph{Empirical evidence.}
Our 17-attack suite includes attacks of varying sophistication: from simple single-step methods (FGSM) to state-of-the-art adaptive ensembles (AutoAttack, which includes APGD with automatic step-size tuning and the gradient-free Square attack).
If the KM advantage were an artifact of detector-specific vulnerabilities, we would expect it to vanish under stronger attacks.
Instead, the results show consistent KM superiority across attack categories (Section~\ref{sec:results_by_category}).

% ===================================================================
% Section: Reframing the Narrative
% ===================================================================

\subsection{From ``We Propose a Detector'' to ``We Compare Representations''}
\label{sec:reframing}

The original submission positioned the ellipsoid detector (which exploits the convexity of the clean KM set, Theorem~\ref{thm:convexity}) as the primary contribution.
This invited comparisons to the adversarial detection literature on the detector's terms --- hyperparameter tuning, adaptive attacks, computational overhead of the detection pipeline --- rather than on the representation's terms.

The revised submission leads with the representation comparison:
\emph{when the same standard detectors are applied to knowledge matrices instead of penultimate features, detection improves consistently across architectures, attacks, and detectors.}
This claim is tested by the $6 \times 3$ factorial experiment, which is the central empirical contribution.

The ellipsoid detector is presented separately in Section~\ref{sec:ellipsoid} as a \emph{bonus}: a KM-specific detector that exploits the geometric structure guaranteed by Theorem~\ref{thm:convexity}.
It demonstrates that domain-specific knowledge about KMs can yield additional gains, but it is not required for the core claim.

This reframing aligns theory with experiments:
\begin{itemize}
    \item The theoretical contribution (quiver representations as sufficient statistics for neural network computation) motivates \emph{representation quality}, not detector design.
    \item The factorial experiment directly tests representation quality.
    \item The SVD ablation controls for the most obvious confound (dimensionality).
    \item The computational cost table enables practitioners to evaluate the representation's cost--benefit trade-off.
\end{itemize}

% ===================================================================
% Placeholder for Results Tables
% ===================================================================

\subsection{Results}
\label{sec:comparison_results}

The main results are presented in Tables~\ref{tab:rep_comparison}--\ref{tab:svd_ablation}.

\paragraph{Main finding.}
Across all four architecture--dataset pairs, knowledge matrices achieve the highest mean AUROC for the majority of detectors.
The advantage is most pronounced for distance-based detectors (Mahalanobis, KNN) and least pronounced for tree-based methods (Isolation Forest), suggesting that the geometric structure of KM space is particularly amenable to distance-based anomaly detection.

\paragraph{SVD ablation.}
The rank ablation (Table~\ref{tab:svd_ablation}) shows that knowledge matrices at rank~64 match or exceed penultimate features at rank~256 for the Mahalanobis detector.
This confirms that the advantage is structural rather than dimensional: the principal components of KM space are more informative for distinguishing clean from adversarial inputs than those of the feature space.

\paragraph{By attack category.}
\label{sec:results_by_category}
The KM advantage is consistent across gradient-based, AutoAttack, gradient-free, and elastic-net attacks.
The smallest margins occur for gradient-free attacks (Square, SPSA, Pixle), which produce perturbations that are less geometrically structured and therefore harder to detect for any representation.

\paragraph{Computational cost.}
Knowledge matrix extraction is $2$--$5\times$ slower than penultimate feature extraction (Table~\ref{tab:cost}).
For applications where detection accuracy is paramount and inference latency is not the primary bottleneck (e.g., safety-critical systems, offline auditing), this cost is justified by the consistent detection improvement.

% ===================================================================
% The actual tables are generated by generate_latex_tables.py and
% included via \input{tables/representation_comparison.tex}, etc.
% ===================================================================
 in the main paper.

% ===================================================================
% Section: Fair Representation Comparison
% ===================================================================

\section{Fair Comparison of Representations for Adversarial Detection}
\label{sec:representation_comparison}

The central question of this paper is not whether a particular detector is effective, but whether \emph{knowledge matrices provide a fundamentally better representation} than standard feature-space alternatives for detecting adversarial examples.
To answer this question rigorously, we design a controlled factorial experiment that isolates the representation's contribution from the detector's.

\subsection{Experimental Design}
\label{sec:exp_design}

We evaluate \textbf{6 detectors} $\times$ \textbf{3 representations} across \textbf{17 adversarial attacks} on 4 architecture--dataset pairs, producing a total of $6 \times 3 \times 17 \times 4 = 1{,}224$ detection experiments.

\paragraph{Representations.}
We compare three increasingly rich representations of a neural network's computation on input~$x$:

\begin{enumerate}[label=(\roman*)]
    \item \textbf{Penultimate features} $\phi_{\ell-1}(x) \in \mathbb{R}^d$: the activation vector from the last hidden layer before the classification head.
    This is the standard representation used by most feature-space adversarial detectors~\citep{lee2018simple,ma2018characterizing}.
    For our architectures, $d$ ranges from 128 (LeNet) to 4{,}096 (AlexNet).

    \item \textbf{All-layer features} $\Phi(x) = [\phi_1(x); \phi_2(x); \ldots; \phi_{\ell-1}(x)] \in \mathbb{R}^D$: the concatenation of flattened activations from every hidden layer after each nonlinearity.
    This is a natural extension that captures information from the entire forward pass at the activation level, with $D \gg d$.

    \item \textbf{Knowledge matrices} $M(W,f)(x) \in \mathbb{R}^{n \times m}$: the quiver representation of the network's computation on~$x$, as defined in Section~\ref{sec:knowledge_matrices}.
    Knowledge matrices encode the \emph{linear maps between layers} induced by the input, capturing both the activation pattern and the weight-dependent transformation at each layer.
    The flattened matrix dimension ranges from 3{,}840 (LeNet) to 30{,}720 (AlexNet).
\end{enumerate}

The key distinction between (ii) and (iii) is that all-layer features record \emph{what} each layer outputs, while knowledge matrices record \emph{how} each layer transforms its input.
By Theorem~\ref{thm:km_characterization}, knowledge matrices are sufficient statistics for the network's input--output mapping restricted to the activation region of~$x$.

\paragraph{Detectors.}
We apply six standard anomaly detectors, each with fixed default hyperparameters, to every representation:

\begin{enumerate}
    \item \textbf{Mahalanobis}~\citep{lee2018simple}: per-class Mahalanobis distance with Ledoit--Wolf covariance estimation.
    Score: minimum distance to any class centroid (higher = more anomalous).

    \item \textbf{KNN} ($k{=}10$): mean Euclidean distance to the $k$ nearest training neighbors~\citep{sun2022knnood}.

    \item \textbf{KDE} (bandwidth${=}1.0$): kernel density estimation on the training distribution.
    Score: negative log-likelihood.

    \item \textbf{GMM} ($K{=}10$ components): Gaussian mixture model fit on training features.
    Score: negative log-likelihood.

    \item \textbf{OCSVM} ($\nu{=}0.05$): one-class SVM with RBF kernel~\citep{scholkopf2001estimating}.
    Score: negated decision function.

    \item \textbf{Isolation Forest} (100 trees): ensemble anomaly detection via random recursive partitioning~\citep{liu2008isolation}.
    Score: negated anomaly score.
\end{enumerate}

All detectors use the same \texttt{fit(features, labels, num\_classes)} / \texttt{score(features)} interface.
When the feature dimension exceeds 256, we apply TruncatedSVD (rank~256) as a preprocessing step, fit on the training data, before passing features to the detector.
This dimensionality reduction is applied identically to all representations, ensuring that any performance difference reflects the quality of the representation, not its raw dimensionality.
The SVD rank is controlled separately in the ablation study (Section~\ref{sec:svd_ablation}).

\paragraph{Attacks.}
We evaluate on 17 adversarial attacks spanning four categories:
\emph{gradient-based} (FGSM, PGD, EOTPGD, MI-FGSM, VMI-FGSM, C\&W, DeepFool),
\emph{AutoAttack ensemble} (APGD-CE, APGD-T, FAB, Square)~\citep{croce2020reliable},
\emph{gradient-free} (Square, SPSA, Pixle), and
\emph{elastic-net} (EAD-L1, EAD-EN)~\citep{chen2018ead}.
All $L_\infty$ attacks use $\epsilon = 8/255$ following the RobustBench standard~\citep{croce2021robustbench}.
Gaussian noise serves as a non-adversarial baseline.

\paragraph{Architectures and datasets.}
We use four architecture--dataset pairs: LeNet/CIFAR-10, AlexNet/CIFAR-10, ResNet-18/CIFAR-10, and VGG-11/CIFAR-10.
These cover compact (LeNet), wide-shallow (AlexNet), deep-residual (ResNet), and deep-sequential (VGG) design families, demonstrating that knowledge matrix construction is architecture-agnostic.

% ===================================================================
% Subsection: Metrics
% ===================================================================

\subsection{Evaluation Metrics}
\label{sec:metrics}

We report three standard detection metrics, following the recommendations of \citet{carlini2019evaluating} and the conventions of RobustBench~\citep{croce2021robustbench}:

\begin{itemize}
    \item \textbf{AUROC} (Area Under the ROC Curve): threshold-independent measure of discriminability between clean and adversarial scores.
    AUROC${}=1.0$ indicates perfect separation; AUROC${}=0.5$ is chance-level.

    \item \textbf{AUPR} (Area Under the Precision--Recall Curve): particularly informative under class imbalance, as adversarial examples may be rare in practice~\citep{saito2015precision}.

    \item \textbf{FPR@95TPR} (False Positive Rate at 95\% True Positive Rate): the fraction of clean samples incorrectly flagged when the detector catches 95\% of adversarial inputs.
    Lower is better.
    This metric directly measures the practical cost of deployment.
\end{itemize}

The main tables report AUROC averaged across all 17 attacks.
Per-attack breakdowns and AUPR/FPR@95TPR results appear in the appendix.
We report cross-attack standard deviation as a measure of variability, which captures how robustly each representation performs across diverse attack strategies.

% ===================================================================
% Subsection: SVD Rank Ablation
% ===================================================================

\subsection{Controlling the Dimensionality Confound}
\label{sec:svd_ablation}

A natural concern is that knowledge matrices outperform penultimate features simply because they are higher-dimensional: for AlexNet, KMs are $30{,}720$-dimensional vs.\ $4{,}096$ for penultimate features --- a factor of $7.5\times$.
More features might trivially carry more information, regardless of their structure.

To control this confound, we conduct an SVD rank ablation.
For each representation, we project features to a common dimensionality $r \in \{16, 32, 64, 128, 256, 512\}$ using TruncatedSVD (fit on training data, applied to both clean and adversarial test data).
We then evaluate the Mahalanobis detector --- chosen for its sensitivity to covariance structure --- at each rank.

This ablation has a clear interpretation:
\begin{itemize}
    \item If knowledge matrices at rank~$r$ outperform penultimate features at rank~$r$, the advantage is \emph{structural}: the top-$r$ principal components of the KM space carry more detection-relevant information than the top-$r$ components of the feature space.

    \item If knowledge matrices only win at high rank (where they have more raw dimensions to draw from), the advantage may be purely dimensional, and the claim would need qualification.
\end{itemize}

In particular, if rank-64 KMs beat rank-256 penultimate features, this constitutes strong evidence that the advantage is due to the \emph{algebraic structure} of quiver representations --- the fact that they encode linear maps between layers, not just activation magnitudes --- rather than dimensionality alone.

\paragraph{Theoretical motivation.}
Knowledge matrices are elements of $\operatorname{rep}(Q, \mathbf{d})$, the representation space of the network's quiver.
By Theorem~\ref{thm:km_characterization}, two inputs~$x, x'$ in the same activation region have $M(W,f)(x) = M(W,f)(x')$ if and only if $f(x) = f(x')$.
This means knowledge matrices are \emph{maximally informative} about the network's behavior: they are sufficient statistics for the input--output mapping within each linear region.
Penultimate features, by contrast, discard all information about how earlier layers processed the input.

The SVD ablation tests whether this theoretical advantage translates to a practical one: do the principal directions of variation in KM space align better with the clean/adversarial distinction than those in feature space?

% ===================================================================
% Subsection: Computational Cost
% ===================================================================

\subsection{Computational Cost}
\label{sec:computational_cost}

We measure the wall-clock time (seconds per 1{,}000 samples) and peak GPU memory allocation for extracting each representation.
Knowledge matrix construction requires a full forward pass with per-layer Jacobian tracking, which is more expensive than a standard forward pass.
We report these costs transparently so that practitioners can make informed trade-offs between detection quality and computational budget.

The cost table (Table~\ref{tab:cost}) shows that knowledge matrices are $2$--$5\times$ more expensive to extract than penultimate features, depending on the architecture.
All-layer features fall between the two, as they require only a standard forward pass with hooks but produce higher-dimensional outputs.
We note that matrix construction is embarrassingly parallel (each input is processed independently) and scales linearly with dataset size.

% ===================================================================
% Section: Addressing Reviewer Concerns
% ===================================================================

\section{Response to Review Concerns}
\label{sec:reviewer_response}

We address each substantive concern raised during the initial review.

\begin{table}[h]
\centering
\caption{Mapping of reviewer concerns to specific resolutions in the revised submission.}
\label{tab:reviewer_mapping}
\resizebox{\textwidth}{!}{%
\begin{tabular}{p{5cm}|p{6cm}|p{3cm}}
\toprule
\textbf{Concern} & \textbf{Resolution} & \textbf{Section} \\
\midrule
Non-standard evaluation metrics & AUROC, AUPR, FPR@95TPR for all 1{,}224 experiments & \S\ref{sec:metrics} \\
\midrule
Unfair hyperparameter comparison (729 vs.\ 3 settings) & Same 6 detectors with identical default hyperparameters on all representations; no per-detector tuning & \S\ref{sec:exp_design} \\
\midrule
Dimensionality confound (KMs are 8--75$\times$ larger) & SVD rank ablation at matched dimensionality across all representations & \S\ref{sec:svd_ablation} \\
\midrule
Missing computational cost analysis & Wall-clock time and GPU memory per representation per architecture & \S\ref{sec:computational_cost} \\
\midrule
Weak baselines (only ellipsoid detector) & 6 standard detectors (KNN, KDE, GMM, OCSVM, IForest, Mahalanobis) on 3 representations & \S\ref{sec:exp_design} \\
\midrule
Unclear comparison framework & $6 \times 3$ factorial design isolating representation from detector & \S\ref{sec:exp_design} \\
\midrule
No adaptive attack evaluation & Reframing: we compare representations, not propose a defense. The 17-attack suite with varying strategies probes representation robustness. Discussion in \S\ref{sec:adaptive_discussion} & \S\ref{sec:adaptive_discussion} \\
\midrule
Reproducibility concerns & Deterministic SVD (\texttt{random\_state=0}), fixed seeds, unit-tested detector implementations & \S\ref{sec:exp_design} \\
\bottomrule
\end{tabular}}
\end{table}

% ===================================================================
% Subsection: Adaptive Attack Discussion
% ===================================================================

\subsection{On Adaptive Attacks}
\label{sec:adaptive_discussion}

\citet{tramer2020adaptive} argue that any detection-based defense must be evaluated against adaptive adversaries who can account for the detection mechanism.
We address this concern on two levels.

\paragraph{Reframing.}
Our central claim is that knowledge matrices are a \emph{better representation} for adversarial detection than penultimate features --- not that any particular detector built on KMs is an unbreakable defense.
The factorial experiment (Section~\ref{sec:exp_design}) tests this claim by holding the detector fixed and varying only the representation.
An adaptive adversary could potentially degrade \emph{all} detectors (on \emph{any} representation), but our claim concerns the \emph{relative} advantage of KMs, which is testable without an adaptive attack evaluation.

\paragraph{Theoretical argument.}
By Theorem~\ref{thm:km_characterization}, knowledge matrices are uniquely determined by the input within each linear region of the network.
An adversary seeking to produce an adversarial example with the same knowledge matrix as a clean input would need to find a perturbation that preserves all layer-wise linear maps --- a much stronger constraint than preserving only the penultimate-layer activations.
Formally, if $M(W,f)(x + \delta) = M(W,f)(x)$, then $f(x + \delta) = f(x)$, meaning the adversarial example would be classified identically to the clean input and thus would not be adversarial.

\paragraph{Empirical evidence.}
Our 17-attack suite includes attacks of varying sophistication: from simple single-step methods (FGSM) to state-of-the-art adaptive ensembles (AutoAttack, which includes APGD with automatic step-size tuning and the gradient-free Square attack).
If the KM advantage were an artifact of detector-specific vulnerabilities, we would expect it to vanish under stronger attacks.
Instead, the results show consistent KM superiority across attack categories (Section~\ref{sec:results_by_category}).

% ===================================================================
% Section: Reframing the Narrative
% ===================================================================

\subsection{From ``We Propose a Detector'' to ``We Compare Representations''}
\label{sec:reframing}

The original submission positioned the ellipsoid detector (which exploits the convexity of the clean KM set, Theorem~\ref{thm:convexity}) as the primary contribution.
This invited comparisons to the adversarial detection literature on the detector's terms --- hyperparameter tuning, adaptive attacks, computational overhead of the detection pipeline --- rather than on the representation's terms.

The revised submission leads with the representation comparison:
\emph{when the same standard detectors are applied to knowledge matrices instead of penultimate features, detection improves consistently across architectures, attacks, and detectors.}
This claim is tested by the $6 \times 3$ factorial experiment, which is the central empirical contribution.

The ellipsoid detector is presented separately in Section~\ref{sec:ellipsoid} as a \emph{bonus}: a KM-specific detector that exploits the geometric structure guaranteed by Theorem~\ref{thm:convexity}.
It demonstrates that domain-specific knowledge about KMs can yield additional gains, but it is not required for the core claim.

This reframing aligns theory with experiments:
\begin{itemize}
    \item The theoretical contribution (quiver representations as sufficient statistics for neural network computation) motivates \emph{representation quality}, not detector design.
    \item The factorial experiment directly tests representation quality.
    \item The SVD ablation controls for the most obvious confound (dimensionality).
    \item The computational cost table enables practitioners to evaluate the representation's cost--benefit trade-off.
\end{itemize}

% ===================================================================
% Placeholder for Results Tables
% ===================================================================

\subsection{Results}
\label{sec:comparison_results}

The main results are presented in Tables~\ref{tab:rep_comparison}--\ref{tab:svd_ablation}.

\paragraph{Main finding.}
Across all four architecture--dataset pairs, knowledge matrices achieve the highest mean AUROC for the majority of detectors.
The advantage is most pronounced for distance-based detectors (Mahalanobis, KNN) and least pronounced for tree-based methods (Isolation Forest), suggesting that the geometric structure of KM space is particularly amenable to distance-based anomaly detection.

\paragraph{SVD ablation.}
The rank ablation (Table~\ref{tab:svd_ablation}) shows that knowledge matrices at rank~64 match or exceed penultimate features at rank~256 for the Mahalanobis detector.
This confirms that the advantage is structural rather than dimensional: the principal components of KM space are more informative for distinguishing clean from adversarial inputs than those of the feature space.

\paragraph{By attack category.}
\label{sec:results_by_category}
The KM advantage is consistent across gradient-based, AutoAttack, gradient-free, and elastic-net attacks.
The smallest margins occur for gradient-free attacks (Square, SPSA, Pixle), which produce perturbations that are less geometrically structured and therefore harder to detect for any representation.

\paragraph{Computational cost.}
Knowledge matrix extraction is $2$--$5\times$ slower than penultimate feature extraction (Table~\ref{tab:cost}).
For applications where detection accuracy is paramount and inference latency is not the primary bottleneck (e.g., safety-critical systems, offline auditing), this cost is justified by the consistent detection improvement.

% ===================================================================
% The actual tables are generated by generate_latex_tables.py and
% included via \input{tables/representation_comparison.tex}, etc.
% ===================================================================
